% CHƯƠNG 4 - Rubric: Dữ liệu, Công cụ và Kịch bản thực nghiệm
\chapter{Dữ liệu và kịch bản thực nghiệm}
\label{chuong:dulieu}

\section{Phạm vi dữ liệu của báo cáo}

Bài báo gốc thực nghiệm trên ba bộ dữ liệu: Mackey-Glass, CHB-MIT và AES
(mục~\ref{sec:cochetuonglai}). Do giới hạn thời gian và việc tải bộ dữ liệu AES cần tài
khoản/API token Kaggle riêng, báo cáo này chỉ \textbf{tái hiện thực nghiệm trên hai bộ dữ liệu
Mackey-Glass và CHB-MIT}; nội dung liên quan đến AES ở Chương~\ref{chuong:phuongphap} chỉ dựa
trên việc đọc mã nguồn và tài liệu của tác giả ở mức lý thuyết, không có số liệu tự chạy để đối
chiếu, và sẽ không xuất hiện lại trong chương này.

\section{Bộ dữ liệu Mackey-Glass}

Chuỗi Mackey-Glass là dữ liệu \textit{tổng hợp}, sinh bằng cách giải số phương trình vi phân
có trễ \eqref{eq:mackeyglass} (mục~\ref{sec:mackeyglass}) với $\tau = 17$, dùng thư viện Python
\texttt{jitcdde}. Tác giả đã công bố sẵn một bản dữ liệu đã sinh (\texttt{data.pkl}) trong kho
mã nguồn; báo cáo này dùng trực tiếp bản dữ liệu này để đảm bảo so sánh được với số liệu công
bố, không tự sinh lại chuỗi mới. Dữ liệu được chia thành tập huấn luyện, kiểm định (validation)
và kiểm tra (test) theo đúng tỉ lệ trong mã nguồn gốc.

\section{Bộ dữ liệu CHB-MIT}

CHB-MIT (Children's Hospital Boston - Massachusetts Institute of Technology)
\cite{shoeb2009chbmit}, công bố công khai trên PhysioNet \cite{goldberger2000physionet}, là bộ
dữ liệu điện não đồ da đầu (scalp EEG) của các bệnh nhân động kinh nhi khoa, lấy mẫu ở tần số
256Hz, ghi theo chuẩn montage 10-20 với 18-23 kênh (tùy bệnh nhân).

Báo cáo này sử dụng dữ liệu của \textbf{ba bệnh nhân: 1, 3 và 5} (trong số 12 bệnh nhân tác giả
liệt kê mặc định cho \texttt{--patient all}), kèm bốn tệp cấu hình nhỏ do tác giả cung cấp sẵn
trong thư mục dữ liệu: \texttt{seizure\_summary.csv} (thời điểm bắt đầu/kết thúc các cơn động
kinh), \texttt{segmentation.csv} và \texttt{special\_interictal.csv} (quy tắc chia đoạn tín
hiệu), và \texttt{sampling\_CHBMIT.csv} (số lượng mẫu ictal/interictal mỗi bệnh nhân, dùng để
tính tỉ lệ oversampling).

\section{Công cụ thực hiện}

Toàn bộ thực nghiệm dùng Python và PyTorch, dựa trên mã nguồn công khai của tác giả (kho
\texttt{SkyeGunasekaran/FutureGuidedLearning}, commit \texttt{5fa24ea}). Mã nguồn đã vá lỗi
được lưu tại kho cá nhân \url{https://github.com/kkyaha/Forecast}, trong đó
\texttt{repro/upstream/} giữ nguyên mã gốc của tác giả để đối chiếu, còn \texttt{repro/build/}
là bản đã vá để chạy được (xem các lỗi phát hiện ở Chương~\ref{chuong:taihien}). Thực nghiệm
Mackey-Glass và Giai đoạn 1 của CHB-MIT được chạy trên máy cá nhân (chip Apple M4, 16GB RAM,
dùng backend MPS của PyTorch); do máy cá nhân hết bộ nhớ khi xử lý dữ liệu EEG cho Giai đoạn
2-3 của CHB-MIT, phần này được chuyển sang chạy trên GPU T4/P100 miễn phí của Kaggle
(Chương~\ref{chuong:taihien}).

\section{Kịch bản thực nghiệm}

\subsection{Mackey-Glass}

Kịch bản chuẩn chạy lệnh
\texttt{exp/base\_exp.py --horizon 5 --num\_bins 50 --alpha 0.5 --epochs 20}, tái hiện đúng
tham số mẫu trong tài liệu hướng dẫn của tác giả, so sánh MSE của Teacher, Baseline và Student
trên tập kiểm tra. Để kiểm tra độ ổn định của kết quả, kịch bản này được lặp lại trên sáu seed
khởi tạo khác nhau ($0, 1, 2, 3, 4, 42$). Ngoài ra, kịch bản \texttt{exp/drift\_exp.py} (mở
rộng của tác giả, kết hợp FGL với phát hiện trôi khái niệm theo kiểm định Page-Hinkley) được
chạy với tầm dự báo quét từ 2 đến 15 bước.

\subsection{CHB-MIT}

Giai đoạn 1 (Teacher, bài toán seizure detection) chạy lệnh
\texttt{exp/seizure\_detection.py --patient \{1,3,5\} --epochs 50 --patience 5}, đánh giá bằng
AUC-ROC. Giai đoạn 2 (Baseline) và Giai đoạn 3 (Student/FGL, bài toán seizure prediction) chạy
\texttt{exp/seizure\_prediction.py} và \texttt{exp/FGL\_CHBMIT.py --patient \{1,3,5\} --epochs
30 --trials 3 --optimizer Adam --alpha 0.5 --temperature 4}. Cả ba giai đoạn chạy trên Kaggle,
từ bản sạch của repo gốc tác giả (ban đầu chỉ bệnh nhân 1 tải được dữ liệu; sau đó tải thêm
được bệnh nhân 3 và 5). Bệnh nhân 1 được chạy lại ba lần độc lập (Giai đoạn 1 còn để đối chiếu
với lần chạy bất thường trên máy cá nhân); bệnh nhân 3 và 5 mỗi bệnh nhân chạy một lần (ba
trial mỗi lần cho Giai đoạn 2-3). Kết quả trình bày trong Chương~\ref{chuong:taihien}.

\subsection{Đối chiếu với số liệu công bố}

Kết quả tái hiện được đối chiếu với hai con số chính mà tác giả công bố: mức giảm MSE 23,4\%
trên Mackey-Glass (loại outlier) và mức tăng AUC-ROC trung bình 44,8\% trên CHB-MIT, cả hai so
sánh mô hình có FGL (Student) với mô hình không có FGL (Baseline). Chi tiết kết quả và mức độ
khớp được trình bày trong Chương~\ref{chuong:taihien} và Chương~\ref{chuong:sosanh}.
